%% =====================================================================
%% Paper-bridge fragment: Abstract + Keywords + Introduction.
%% Replaces the abstract/keywords block (Main.tex L72--82, class
%% frontmatter idiom preserved: \begin{abstract} ... \keyFont{...}) and
%% the Introduction (Main.tex L84--95). Body content only; no preamble.
%% =====================================================================

\begin{abstract}
\section{}
Artificial muscles such as braided pneumatic actuators (BPAs) reproduce the force and contraction behavior of skeletal muscle, making them attractive for the legs of musculoskeletal robots. However, the isometric torque a BPA delivers about a joint remains difficult to predict, because rigid-body torque models omit attachment compliance, effective-length offsets, and the contraction lost where the actuator wraps over the joint. In this work, we measured isometric knee torque on pinned and biomimetic knee test stands across actuator resting lengths, isolated error sources with a hybrid torque calculation, and identified correction terms by multiobjective optimization. Our analysis resulted in (1) a corrected torque model whose length offset, effective stiffness, and fictive-strain wrapping loss, expressed through the unloaded \qty{620}{\kPa} contraction, halve prediction error on held-out tests across configurations; (2) a cross-configuration-consistent stiffness pair selected by a simultaneous design requirement; and (3) redesigned attachment routes that meet or exceed human isometric knee-torque targets with quantified margins, confirmed by first hardware measurements. The corrected model and redesign pipeline support first-iteration torque design of BPA-actuated robot joints.

\tiny
 \keyFont{ \section{Keywords:}
Braided Pneumatic Actuator, Artificial Muscle, Isometric Torque, Moment Arm, System Identification, Actuator Routing, Biomimetic Robotics, Knee Joint}
\end{abstract}

\section{Introduction}

Musculoskeletal robots built on biological design principles are experimental platforms for motor-control questions that are impractical or unethical to study in human subjects, from spinal reflex loops to the coupling of posture and locomotion \citep{ijspeert_amphibious_2020, goldsmith_investigating_2021, liu_robotic_2018}. The same principles also yield machines that copy the anatomy of their biological counterparts: human-mimetic humanoids and anthropomorphic arms with anatomically arranged musculotendon networks \citep{asano_design_2017, asano_musculoskeletal_2019, hitzmann_anthropomorphic_2018}, lower-limb robots driven by multifilament muscles \citep{kurumaya_musculoskeletal_2016}, and platforms that isolate the function of a single muscle such as the gluteus medius \citep{shin_constructive_2018} or select actions during adaptive locomotion \citep{schilling_adaptive_2022}. Each of these machines asks its actuators to do what skeletal muscle does: pull as a lightweight tension member whose force depends on excitation and length, routed from origin to insertion across a joint. The joint-level quantity that matters is not force but torque, and torque depends as much on the routed path as on the actuator itself.

No actuator family reproduces muscle exactly, and each alternative trades something away. Geared electric motors dominate position control, but they place mass and reflected inertia at the joint and gain compliance only through added series elasticity \citep{veale_towards_2016}. Hydraulics supply high force density at the price of pumps, valves, and rigid supply lines \citep{suzumori_trends_2018}. Shape-memory alloys respond slowly, and dielectric elastomers deliver low force at impractical drive voltages \citep{madden_artificial_2004}. Braided pneumatic actuators (BPAs), also called McKibben artificial muscles, occupy the practical middle ground: a pressurized sleeved bladder that shortens and pulls, producing force that falls with contraction in a muscle-like way while remaining a compliant tension member. Air-muscle quadrupeds have demonstrated robust, adaptive gaits \citep{aschenbeck_design_2006}, and the Festo fluidic BPAs used in this work are commercially available in the sizes a human-scale leg requires.

For BPAs, the force side of the design problem is well characterized: we have published closed-form equations for the maximum force and the isometric force of \qtylist{10;20}{\mm} Festo actuators as functions of pressure, contraction, and resting length \citep{bolen_2026}. The torque side has received far less attention. Joint torque is the product of actuator force and the moment arm of the actuator's line of action about the joint axis, and the moment arm is a property of the attachment route that changes with joint angle. Biomechanics treats moment-arm determination as its own measurement problem, from lower-limb muscular anatomy \citep{brand_model_1982} and three-dimensional musculotendon estimation \citep{menegaldo_moment_2004} to curved paths with moving attachment points \citep{suderman_moving_2012}, because short moment arms amplify path error \citep{young_analyzing_2019}. A robot designer adds effects that no biological model carries: brackets and fixtures deform under load, artificial tendons and cabling stretch, and at deep flexion the actuator bends over the joint and loses usable contraction. None of these effects appears in a rigid-body torque model built from a characterized actuator and nominal geometry.

Torque modeling for braided actuators has consequently concentrated on the force law and on placement. \citet{chou_measurement_1996} derived the static force relation of McKibben muscles from braid geometry and quantified its hysteresis, and later work refined static models and approximations for pneumatic and fluidic muscles \citep{tondu_static_2022, sarosi_new_2012}. For the actuator family used here, \citet{martens_modeling_2017} modeled static force from a geometric and material framework, and \citet{hunt_modeling_2017} modeled length effects. At the joint level, \citet{bolen_determination_2019} determined artificial-muscle placement for biomimetic humanoid robot legs, and \citet{morrow_optimization_2020} optimized placements so that a bipedal robot matches human torque and range-of-motion capabilities. What these models do not provide is a torque prediction that survives contact with a physical joint. None of them represents fixture and tendon compliance in the load path, carries an effective-length offset between modeled and physical muscle length, or accounts for the contraction lost where the actuator wraps over the joint. To our knowledge, no BPA torque model has been validated across joint configurations, actuator resting lengths, and both antagonistic sides of a joint with a single consistent parameter set.

The present work treats these mispredictions as identifiable design factors, not residual noise. We measured isometric knee torque on two test stands: a robot leg with a pinned revolute knee, and the same leg with a biomimetic knee whose instantaneous center of rotation migrates with flexion like that of the human joint \citep{steele_biomimetic_2018, steele_development_2017}. Each knee was driven by antagonistic flexor and extensor BPAs of \qtylist{10;20}{\mm} diameter across a range of resting lengths, with and without artificial tendons. A hybrid torque calculation, which combines measured muscle lengths and moment arms with modeled actuator force, isolated kinematic error from force error before any parameter was fitted.

Our analysis resulted in (1) a corrected torque model whose constant length offset ($\chi_0$), lumped attachment stiffness ($\chi_1$ and $\chi_2$), and fictive-strain wrapping loss ($\chi_3$) halve prediction error on held-out tests across configurations; (2) a two-stage multiobjective identification that holds one stiffness pair consistent across both knees and both actuators, validates it on held-out resting lengths and on the biomimetic knee, and selects the adopted set by an external design requirement; and (3) redesigned attachment routes, each driven by two \qty{20}{\mm} BPAs, that meet or exceed human isometric knee-torque targets from OpenSim musculoskeletal models \citep{seth_opensim_2018, delp_interactive_1990} with a required \qty{5}{\percent} margin at every angle, together with the first hardware torque measurements of the redesigned extensor route. Throughout, the wrapping loss is expressed through fictive strain: the excess contraction the force model must see beyond the actuator's physical contraction, normalized by the unloaded \qty{620}{\kPa} contraction $\varepsilon_{620}$ of the free actuator. Together, the corrected model and the redesign pipeline let a designer predict isometric joint torque, correct it against hardware, and route the actuators to a human-derived target on the first design iteration.
